% answers_short.tex —— 2011 年·江西卷·文科 · 简短答案速查
% 由 tools/make_short_answers.py 自 answers.tex 的 \gksolhead/\gkanswer 提取，
% 并与 2011/2011_数学真题_江西文_答案要点.md 逐题比对后生成，勿手改（改源 MD 后重跑）。

\begin{gknotice}
\textbf{使用说明}：本节只给最终答案便于快速核对，完整推导见解析版。答案为编写者独立推导，\textbf{非官方标准答案}。
\end{gknotice}

\gkpart{选择题}
\begin{center}\begin{tabular}{ccccccccccc}
\toprule
题号 & 1 & 2 & 3 & 4 & 5 & 6 & 7 & 8 & 9 & 10 \\
\midrule
答案 & B & D & C & A & B & B & D & C & D & B \\
\bottomrule
\end{tabular}\end{center}

\gkpart{填空题}
\begin{itemize}[leftmargin=2.6em]
  \item[\textbf{11.}] $-6$。
  \item[\textbf{12.}] $48$。
  \item[\textbf{13.}] $6$。
  \item[\textbf{14.}] $-8$。
  \item[\textbf{15.}] $[0,+\infty)$。
\end{itemize}

\gkpart{解答题}
\begin{itemize}[leftmargin=2.6em]
  \item[\textbf{16.}] (1) $\dfrac1{10}$；(2) $\dfrac7{10}$
  \item[\textbf{17.}] (1) $\dfrac13$；(2) $\dfrac{\sqrt3}2$
  \item[\textbf{18.}] (1) $\dfrac{2\sqrt3}3$；(2) 证明见解析
  \item[\textbf{19.}] (1) $y^2=8x$；(2) $2$
  \item[\textbf{20.}] (1) $f(x)=\dfrac13x^3+3x^2+2x$；(2) $(m,n)=(2,3)$ 或 $(3,5)$
  \item[\textbf{21.}] (1) $0<a\leqslant\dfrac13$；(2) 不存在
\end{itemize}
